\chapter{Overview}
\section{Purpose}
This guide explains how EFC approaches the design, implementation, testing, commissioning, documentation, operation, and maintenance of industrial control systems. It is written primarily for engineers who are new to EFC projects and is intended to help them build systems that are safe, supportable, repeatable, secure, and consistent.

\section{Scope}
The guide covers PLCs, remote I/O, industrial control panels, operator workstations, HMI/SCADA systems, industrial networks, instrumentation interfaces, packaged-equipment integration, and the related project documents. It applies to internal engineering work and to contractors unless a project specification gives a more stringent requirement.

This document does not replace the adopted electrical code, the requirements of the Authority Having Jurisdiction (AHJ), a project risk assessment, a process hazard analysis, or a machine-safety/functional-safety assessment. The project engineer must identify the governing jurisdiction and applicable editions at project initiation.

\section{How to Use This Guide}\label{sec:how-to-use}
You don't need to read this guide cover to cover. Start with the parts that match where your project is, and come back to the rest when you need it.

\subsection*{Where to start}
\begin{itemize}[itemsep=2pt]
  \item \textbf{New to EFC?} Read Chapter~\ref{ch:terminology} (Terminology) and Chapter~\ref{ch:design-process} (Design Process) first. Chapter~\ref{ch:design-process} walks through every project stage and tells you what you need going in, what you do, and what you hand off.
  \item \textbf{Picking parts or designing a panel?} Go to Chapter~\ref{ch:hardware} (Hardware). It lists our standard parts and explains the panel design checks (SCCR, heat, wiring).
  \item \textbf{Writing PLC or SCADA code?} Use the PLC and SCADA standards chapters and the appendices.
  \item \textbf{Working with a contractor or panel shop?} See the Working with Contractors chapter for what they need to send us.
\end{itemize}

\subsection*{What the key words mean}
\begin{itemize}[itemsep=2pt]
  \item \textbf{Must} --- required. It's needed for safety, code compliance, security, or a basic EFC practice. If you can't meet it, talk to the Controls Engineering Manager before you go further.
  \item \textbf{Should} --- the normal way we do it. You can do something different if the project needs it; just write down why.
  \item \textbf{May} --- an option you're allowed to use.
\end{itemize}

\subsection*{How to read the notes in this guide}
\begin{itemize}[itemsep=2pt]
  \item \textit{\textcolor{noteorange}{[Orange bracketed text]}} marks an open question or something that still needs a decision. Don't treat it as settled.
  \item \textit{\textcolor{futureblue}{[Future: blue bracketed text]}} marks something we plan to add later. It is not part of the current process.
  \item A green \newtag\ tag marks content that was recently added.
  \item Blue text is a link. Links to manufacturers and outside references are for convenience; keep a controlled copy of any datasheet the project depends on.
\end{itemize}

\subsection*{When the guide doesn't fit your project}
Every project is a little different. If something here doesn't fit, or you think it's wrong, raise it with the Controls Engineering Manager. Changes for one project are recorded as a deviation (see Section~\ref{sec:governance}); changes that should apply to everyone go into the next revision of this guide.

\section{Governance}\label{sec:governance}
Controls Engineering owns this guide, and the Controls Engineering Manager (or someone they choose) approves changes to it. We review it at least once a year, and sooner if a code, a main platform, a security requirement, or a standard part changes in a big way.

If a project needs to do something differently, write it down in the design basis as a deviation. Include what requirement you're not meeting, why, what the risk is, what you're doing instead, who approved it, and when it will be closed out. A deviation only applies to that one project; it doesn't change the guide for anyone else.

\section{Controlled Layers}
Controls documentation must be maintained in four layers:
\begin{enumerate}[itemsep=2pt]
  \item Company-wide mandatory engineering requirements.
  \item Approved hardware and software catalog, including lifecycle status.
  \item Version-controlled reusable PLC, SCADA, and drawing libraries.
  \item Project-specific design, test, configuration, and validation records.
\end{enumerate}
